\documentclass[a4paper,14pt]{extarticle}
\usepackage{graphicx}

% Пакеты для русского языка и кодировки
\usepackage[T2A]{fontenc}
\usepackage[utf8]{inputenc}
\usepackage[russian]{babel}

% Пакеты для математики и теорем
\usepackage{amsmath, amssymb, amsthm}

% Пакеты для графики и гиперссылок
\usepackage{hyperref}

% Настройка полей
\usepackage{geometry}
\geometry{top=2cm, bottom=2cm, left=2.5cm, right=1.5cm}

% Для оформления кода и таблиц
\usepackage{listings}
\usepackage{longtable,booktabs,calc}
\usepackage[table]{xcolor}

% Настройка межстрочного интервала
\usepackage{setspace}
\onehalfspacing

% Отключает нумерацию для part, chapter, section и т.д.
\setcounter{secnumdepth}{-3}

\hypersetup{
    colorlinks=true,
    linkcolor=black,
    urlcolor=blue
}


\begin{document}
\thispagestyle{empty}

\begin{center}
    {\Large\textbf{Итоговая оценка за курс и экзамен}}
\end{center}

\section*{Формула оценки}

Итоговая оценка складывается из оценок за два раздела курса:

\[
\begin{aligned}
    \text{О}_{\text{курс}} &= 0{,}5\,\text{О}_1 + 0{,}5\,\text{О}_2,\\
    \text{О}_1 &= \text{КР}_1,\\
    \text{О}_2 &= 0{,}25\,\text{ЛБ}_4 + 0{,}25\,\text{ЛБ}_5
        + 0{,}10\,\text{Активность} + 0{,}40\,\text{КР}_2.
\end{aligned}
\]

Рассчитанная оценка за курс округляется до целого по арифметическому правилу:
дробная часть от 0{,}5 округляется вверх. \textbf{Округлённая накопленная
оценка ниже 4 означает, что курс не сдан.}

\section*{Последнее занятие и контрольная работа №2}

Шестое занятие, которое пройдёт \textbf{15--21 октября}, --- последнее. На нём
состоятся контрольная работа №2 и защита лабораторной работы №5. В контрольную
работу №2 войдут темы лекций №4--6 и лабораторных работ №4--5.

\section*{Экзамен}

Повысить итоговую оценку можно на письменном экзамене. Студенты, чья округлённая
накопленная оценка \textbf{ниже 4}, записываются на него автоматически.

Студентам с накопленной оценкой \textbf{4 или выше}, которые хотят сдавать
экзамен, нужно написать преподавателю.

\textbf{Оценка за экзамен полностью заменяет накопленную оценку за курс.} Если
записавшийся студент не придёт на экзамен, его оценка за экзамен и итоговая
оценка за курс будут равны \textbf{0}.

Экзамен будет письменным и охватит все темы лекций, семинаров и лабораторных
работ.

\end{document}
